\section{Introduction}
\label{sec:main-introduction}

A passive network does not choose its route to optimize a transport cost.
Conservation and the potential drops across its edges determine the
route. This distinction matters when nominations or physical coefficients
are uncertain: even though each scenario has a unique state, finding a
scenario with the largest pressure difference or arc flow is a nonlinear
global optimization problem. It is also an operationally useful problem.
A worst-case arc flow tests a capacity, a potential difference measures
a required pressure range, and a weighted potential objective measures
several locations together. A tractability statement for one of these
questions need not answer the others.

Potential-based models arise in passive gas and water transport and in
nonlinear resistive circuits. The survey of
\citet{pfetsch2026-potential-based-flowsan-overview} describes their common
structure and their different physical interpretations. We work with the
abstract passive equations, principally with the quadratic law
$\pi_u-\pi_v=\beta_e x_e|x_e|$. The potential variable may already be a
physical transformation, such as squared pressure. Our results concern
this stated mathematical model; they do not replace a system's complete
engineering model, controls, or admissible operating conditions.

Four choices govern the complexity developed here. First, cycles may
be independent blocks or interact inside a block. Second, an objective
may inspect one arc, compare two potentials, or sum measurements across
the graph. Third, uncertain coefficients may vary independently in
intervals, belong to finite sets, or satisfy shared constraints. Fourth,
the requested output may be an exact threshold decision or an additive
interval and an admissible near-optimal scenario. Each distinction leads
to a provable boundary, often on a very small graph class.

The last distinction is particularly easy to overlook. Two parallel
quadratic paths already produce a square root. A chain of independent
cycles accumulates many such roots. Separately computing every local
state exactly does not give an efficient exact comparison of their sum.
Yet each local contribution can be enclosed to enough bits that the
sum is accurately known. We therefore measure additive running time in
$\log(1/\eps)$, not in $1/\eps$, and require the returned rational
parameters to satisfy the original scenario constraints exactly. Their
physical state can remain irrational.

\subsection{Contributions and their scope}

The main results form three connected groups. The first group concerns
nomination uncertainty and arithmetic. A law-independent face reduction
localizes the free nominations of a terminal-drop maximizer to one
block and bounds their number there. Independent coefficient choices
and objective contributions may still span the whole block path. Combined
with fixed-dimensional real algebraic computation, it yields polynomial
additive optimization at each fixed maximum block rank, including joint
independent coefficient uncertainty and growing dense polynomial laws.
An individual arc is local to one block and therefore admits an exact
algebraic answer in the polynomial-law model. On cacti, exact
single-source/sink pressure comparison is instead polynomial-time
equivalent to Square-Root Sum. Fixed fractional powers preserve the
additive algorithms but can make exact arc comparison Square-Root-Sum
hard on a single cycle.

The second group identifies the role of the observable. Under fixed
nominations, the universal graph class on which resistance sets may be
replaced by interval hulls is exactly the trees for weighted potentials,
the cacti for pairwise potential differences and linear flows, and the
series-parallel graphs for individual arcs. The attainable flow region
is universally convex exactly on cacti; normalized potential and joint
regions are universally convex exactly on trees. These geometric facts
explain both endpoint algorithms and small-graph counterexamples. They
do not by themselves remove exact-arithmetic or nomination hardness.

For weighted potential optimization, a stronger face theorem leaves
$O(r_0p)$ free nomination coordinates when $r_{\max}\le r_0$ for a
positive bound $r_0\ge1$ and objective support is at most $p$. We then prove an additive
compiler for arbitrarily many quadratic cactus blocks at fixed affine
nomination dimension, allowing independent directional coefficients.
A scalar semialgebraic approximation construction extends the method to
arbitrarily many bounded-rank blocks on one nomination line, including
fixed dense piecewise-polynomial laws. Several nomination parameters
are allowed when the \emph{total} rank in noncactus blocks is fixed.
These are completed extensions. The remaining unrestricted sum of
higher-block responses in several shared parameters is stated separately
in \cref{eq:a-wblk-residual}.

The third group concerns continuous design and verifiable computation.
For fixed nominations on a cactus, arbitrary-dimensional parameter
polytopes within individual cycles preserve scalar-root structure.
We give exact rational parameter recovery for arc extrema and local
capacities, additive minimization of dense convex polynomial flow
objectives, and additive convex maximization through a fixed number of
linear measurements. Coupled convex maximization is hard even with
independent intervals; global coefficient correlations make total-flow
maximization and minimum dissipation strongly hard on cacti. In the
opposite direction, maximum dissipation over a positive rational
resistance polytope admits a polynomial additive algorithm on every
graph. Rational primal-dual, Bregman, support, and curvature witnesses
then certify physical quantities and uncertainty scenarios without
trusting a numerical solver's status.

The inspected primary sources do not supply the combined fixed-block-rank
nomination and coefficient guarantees of \cref{thm:main-block}, the
weighted scalar and hybrid extensions of \cref{thm:main-weighted}, or
the precise objective-dependent classifications and output boundaries
proved here. We identify these scoped results and their network-specific
proofs as contributions of this manuscript, within the limits of the
literature comparison below. The
energy principle, electrical adjoints, circuit confluence, parametric
linear optimization, quantifier elimination, approximation by analytic
panels, and zonotope enumeration are classical ingredients and are
credited where used. We make no general firstness claim about nonlinear
network sensitivity or endpoint evaluation.

\subsection{Relation to earlier work}
\label{subsec:main-literature}

\paragraph{Passive states and circuit structure.}
The primitive-energy characterization goes back to nonlinear network
theory. \citet{duffin1947-nonlinear-networks-iia},
\citet{birkhoff1956-non-linear-network-problems}, and
\citet{minty1960-monotone-networks} establish variational and monotone
network foundations under their respective law hypotheses. We supply
a self-contained existence proof for continuous strictly increasing
resistance laws vanishing at zero, including bounded laws; surjectivity
is unnecessary in this formulation. The grounded Laplacian and adjoint
calculation also have direct predecessors: the conjugate-energy gradient
and Hessian in \citet[Section IV-A]{misra2015-maximum-throughput-problem-in-dissipative}
give potentials and the inverse weighted Laplacian.

The topology-based sign mechanism is classical circuit confluence
\citep{duffin1965-topology-of-series-parallel-networks}.
\citet[Theorem 3 and conclusion]{chauffoureaux1990-monotonicity}
use a linear-circuit criterion to study nonlinear parameter monotonicity
and extremal evaluation. The undated Wang--Hasler report studies
convexity of scalar transfer characteristics as source values vary,
including structural sufficient conditions and ladder networks
\citep[Theorems 1, 4, and 5]{wanghasler1997-convexity}.
Our full-region convexity statements concern a different variable and
object: the vector of attainable flows or potentials under independent
resistance variation. The 1993 paper \emph{Parameter tolerances in non-linear resistive
circuits: worst case analysis based on monotonicity} by Hasler and Wang
\citep{hasler1993-parameter-tolerances} could not be inspected in full, so the comparison does not support an
unrestricted priority claim about endpoint or envelope identities.

\paragraph{Robust feasibility and network complexity.}
The joint product of a balanced nomination box and positive resistance
intervals is established in \citet[Section 3.3]{amann2019-exact-methods-for-two-stage}.
The circulation sign-region method and cycle-capacity linearization
are already present in
\citet{amann2018-deciding-robust-feasibility-and-infeasibility}.
Single-cycle booking validation has a polynomial Turing-time algorithm
under the rational-data hypotheses of
\citet[Section 6]{labbe2021-deciding-feasibility-on-a-cycle}.
Our block-rank algorithms allow an unbounded number of blocks and total
cycles, while handling the arithmetic of their combined objective
explicitly. Th\"urauf's general nomination hardness construction
\citep{thurauf2022-deciding-the-feasibility-of-a} supplies an important
boundary: its main block has unbounded rank although the graph is
series-parallel. We credit that construction and derive the required
bit-gap and arc-probe consequences.

\citet[Lemma 4.2 and Theorem 4.3]{gro2017-algorithmic-results-for-potential-based}
give homogeneous series-parallel reductions and discuss approximate
root arithmetic in the Turing model. Our exact weak comparison results
separate the cost of carrying out those reductions from the cost of
ordering their resulting radical expression. The January 2026 overview
still asks for the complexity of nonlinear cactus maximum potential
difference \citep[p.~10]{pfetsch2026-potential-based-flowsan-overview}.
We resolve its additive version and classify the exact single-source/sink
subclass. We do not claim an exact classification for arbitrary cactus
nomination boxes. In particular, Square-Root-Sum hardness is neither
an NP-hardness result nor evidence of NP membership.

\paragraph{Local formulas, approximation, and design.}
Gotzes, Heitsch, Henrion, and Schultz derive quadratic-radical circulation
formulas on affine nomination rays and discuss node-disjoint cycles
with attached trees \citep[Theorem 6, pp.~18 and 24 of the author manuscript]
{gotzes2016-quantification}. Vigneron combines algebraic functions through
common arrangements and approximate summation; his bit-model schemes
have polynomial dependence on $1/\eps$
\citep[Theorem 6 and Sections 2.3, 3.2]{vigneron2014-geometric-optimization-and-sums-of}.
Our weighted results require polynomial dependence on accuracy bits,
with all surrogates expressed in the same original nomination
coordinates and with rational original-parameter recovery.
\Cref{sec:a-weighted-blocks} compares the scalar construction precisely
with analytic enclosure, rational algebraic approximation, and
semialgebraic parametrization results
\citep{petras2002-approximation,borcea2006-rational-algebraic,yomdin2014-parametrizations,binyamini2019-complex-cells}.
The difference in interface matters; no new general principle of
analytic approximation is asserted.

The exact cycle-root recovery refines a classical LP-bisection mechanism
\citep{agrawal2020-dqcp,megiddo1979-rational} with dense-degree separation
and parameter output bounds. Few-measurement convex maximization uses
zonotope methods \citep{onn2004-convex,ferrez2005-fixed-rank}.
The recent design work of
\citet{klimm2026-approximating-the-network-design-problem} chooses installed
arcs and conductances under costs and a potential range. Its Corollary 4
uses convexity when fixed costs vanish; Theorem 11 gives an FPTAS on
series-parallel networks with strictly positive variable costs and finite
conductance bounds. Those investment objectives and relative-error
contracts differ from optimization over our prescribed resistance
scenario domain. Similarly,
\citet{thurauf2025-adjustable-robust-nonlinear-network-design} reduce
fixed-design robust checking to polynomially many nonlinear worst-case
subproblems and use an adversarial design algorithm; a polynomial number
of subproblems does not itself establish polynomial complexity for each.

\subsection{Reading the paper}

The main text states the principal results, explains their mechanisms,
and illustrates the boundaries. \Cref{sec:main-model} fixes conventions;
\cref{sec:main-localization,sec:main-observables} treat block rank,
arithmetic, and topology; \cref{sec:main-weighted,sec:main-design} develop
weighted objectives and correlated design; and
\cref{sec:main-certificates} explains the certificates and reproducible
evidence. The complete proofs and all supporting results are included
in Appendices \ref{sec:a-pre-preliminaries}--\ref{sec:a-cert},
in dependency order. They are part of this manuscript. No theorem
requires a repository note or an unwritten companion paper.
